\section{Conclusion}
\label{sec:conclusion}

We have shown that GPU-accelerated 5G vRAN is memory-bound, not compute-bound. The dominant kernel in the uplink equalization pipeline---coefficient application---has an arithmetic intensity of 4.32~FLOP/byte, well below the A100's Ridge point of 161~FLOP/byte. Even after Tensor Core optimization yielding a 5.5$\times$ speedup, the kernel achieves only 0.86\% of peak TC throughput, limited by L1 and DRAM bandwidth. This memory-bound regime is structural: it persists across antenna counts, bandwidths, and symbol counts.

Through six controlled experiments on an A100, we established that multi-cell memory contention is governed by L2~cache capacity (42~MB), not HBM bandwidth. The key findings are: (1)~degradation correlates with working set overflow beyond L2, not with kernel type; (2)~pipeline stage ordering by ascending working set reduces contention by up to 52\%; (3)~online contention prediction from 5~latency samples achieves $<$3\% error; (4)~batched serial execution with batch=2 is the optimal latency--throughput tradeoff. The aggregate bandwidth saturates at $\sim$780~GB/s (40\% of HBM peak), providing a calibrated contention model $\delta \approx C/N$.

Based on this evidence, we proposed a memory-hierarchy-aware scheduler combining L2-aware stage co-scheduling, online contention prediction, and adaptive batch sizing with deadline guarantees. The scheduler exploits the vRAN pipeline structure---a feature absent in DNN serving workloads---to reduce contention through stage ordering, a novel observation validated experimentally. The development roadmap toward a full NSDI submission spans 5--7~months, encompassing scheduler prototype, Aerial cuPHY integration, multi-cell deployment, and baseline comparison.
